\section{楚：疆域很大，中心却未必更稳}

江汉、淮河和更南的山川，为楚保留了别国难有的回旋地带。一座城失守以后，王室、兵众和粮道还可能向东或向南移动；北方的敌手却不能轻易用一次会战抹平这样大的政治体。不过，这幅优势不能画成边界清晰、资源会自行流向都城的地图。楚廷要使田租成为军粮、使城邑承担守备、使封君与县邑的命令能够相互辨认，命令便要越过很长的道路，也要不断与宗室、封君和地方官所掌握的关系周旋。楚的问题不是土地够不够，而是谁能够在需要时把这些土地上的人、粮、舟车和文书接到同一条调度线上。

这种困难在楚悼王即位前后的政治选择中已经显露。约前417年，传世国别叙事与后来的楚简研究所提示的巡察、赋役和地方兵粮问题，被归入\eventanchor{WQ-0417-CHU-DAO-INSPECTION}{战国·楚悼王核察地方}的背景；前401年悼王即位，\eventanchor{WQ-0401-CHU-DAO-ASCEND}{战国·楚悼王即位}。这不是说我们保存着一份逐邑核查的王命。较可把握的是，面对三晋的北方压力和国内分散的封君资源，王权需要更多关于地方的消息，也需要能把消息变成军政处置的人。吴起得以进入楚廷，正是在这个缺口中。

\subsection{印信、竹简与广域统治的摩擦}

把楚写成只靠贵族口头支配的南方大国，同样会遮住材料所见的另一面。楚系官玺、封君玺与\eventanchor{WQ-0393-CHU-SEALS}{战国·楚系印信}，说明转运财物、传递命令和确认身份已经需要可辨的凭信；约前四世纪前期，地方官以简牍处置田宅、诉讼和赋役的实践，成为\eventanchor{WQ-0372-CHU-COUNTY-WRITING}{战国·楚地文书行政}所概括的线索。包山墓所出简牍及其司法、封君事务，又以一处保存下来的档案片段，构成\eventanchor{WQ-0316-BAOSHAN}{战国·包山楚简}的物证。它们让人看见，封君权力和文书程序可以并存：印信并没有取消地方关系，简牍也没有使王命无摩擦地抵达每一个邑。

这一点关系到楚的资源为何不自动等于国家能力。对都城而言，清点户口、守备、仓储、田租、渡口和差役，都是把广阔地域转成可用资源的必要环节；对邑吏和编户而言，它们则是更细密的负担。前335年前后楚在淮北、方城方向整修防线，\eventanchor{WQ-0335-CHU-NORTH-DEFENSE}{战国·楚的北方防务}；账本所列同年及前313、前292、前240、前223年前后的户口、仓储、道路、田宅和官吏责任，正可用来说明这种防务会提出哪些行政要求。

但要求不是已经逐条落实的制度。传世《史记》和《资治通鉴》的相关叙述并未留下这些年份的一组完整诏令；出土材料也只是局部。因而，不能把“清核”“重排”“整治”复原成一套每年按期颁行、全楚同样执行的改革。它们更稳妥地标示了战争中的治理工作：中心若要得到粮与兵，就必须设法核验城邑、仓廪、交通和人身关系；地方能否照办、办到何种程度，仍取决于封君、官吏、道路和战局。楚有行政的工具，却不能由此推出它已经拥有后来秦式的均质穿透力。

\begin{sourcenote}[从包山简到王廷政策，不能跳过的距离]
\sourcebook{包山楚简}、楚系玺印和墓葬发掘资料，是接近楚国地方手续的材料，适合说明诉讼、田宅、人身关系与验信并非纯粹口传；它们保存地点有限，不能替全楚作统计。关于前335、前313、前292、前240与前223年前后城邑清核、仓储、道路和差役的安排，传世史书提供的是战局与治理压力的脉络，不能当作逐条政策文本。本文据此区分可见的文书实践、后出材料中归纳出的治理方向，以及尚不能确证的全国实施。
\end{sourcenote}

\subsection{吴起：从魏国军政经验到楚廷的联盟难题}

吴起并非在楚的宗族网络中成长。传世材料把他在魏的用兵、训练与早期军政声望，连到\eventref{WQ-WEI-WUQI}{战国·魏武卒}的传统上。这个“履历”不能由《吴子》中整齐的选拔标准倒推出一份当时的军籍簿；但它说明，楚悼王所起用的不是只会陈说的一位外来游士，而是被后世记忆为熟悉军队、赏罚和国家竞争的人。这样的外来性也有代价：他的权力更依赖悼王的支持，难以自成一套稳固的楚国政治根基。

\begin{event}{约前382年至前381年}{吴起的改革与贵族的反应}{年代与细节有争议}{楚}\eventanchor{WQ-0382-CHU-WUQI-PRELUDE}{战国·吴起改革前奏}\eventanchor{WQ-0382-WUQI-FENGJUN-CUTS}{战国·吴起削封君余荫}\eventanchor{WQ-0381-WUQI-CHU}{战国·吴起在楚变法}
\sourcebook{史记·孙子吴起列传}与《吕氏春秋》将削减封君余荫、约束世禄、整饬官吏和罢去冗员等措施归到吴起与悼王名下。若这些方向能够落到执行，所得便不只是省下若干俸禄：原本附着于世卿、宗室和封邑的资源，才可能较多地转向王室的军费、任官与赏罚。触及的也不只是抽象的“旧制”。被削减的俸禄、世袭的优待和由此连出的地方关系，正是许多精英维持位置的凭借。

不过，传世文本并非改革诏书汇编。《韩非子·和氏》以及后出的叙事能显示后人怎样把楚的难题理解为封君资源与王权军政的争夺，却不足以逐县确认每项措施何时发布、由谁执行、持续多久。应当把吴起的改革看作一场有明确方向的政治重组，而不是一套条文齐备、无阻运行的现代行政方案。它需要悼王、执行官与部分地方力量共同承受成本；没有这个联盟，政策的文字或倡议本身不会自动变成国家能力。
\end{event}

悼王死去时，改革的支点随之消失。前381年楚肃王继位，\eventanchor{WQ-0381-CHU-SU-ACCESSION}{战国·楚肃王继位}；同一权力交接为受损贵族重新进入朝政打开了窗口。吴起在丧礼中遇害的基本结局，见于\eventanchor{WQ-0381-WUQI-DEATH}{战国·吴起被杀}；《史记》所写他伏在悼王遗体旁、反对者射中遗体、肃王后来诛及相关家族的经过，则由\eventanchor{WQ-0381-CHU-WUQI-CORPSE}{战国·吴起丧礼遇害}保留为格外有力的政治仪式叙事。

\begin{debate}[丧礼场景说得有多实]
吴起被杀、其死与贵族反弹有关，是《史记·孙子吴起列传》与《吕氏春秋·贵卒》共有的叙事核心；伏尸、射尸和事后诛族的连锁细节，更适合被读作后出史家用来说明王体不可侵犯、继承时刻危险性的场景组织。它不能支持虚构现场对话或揣测每个行动者的内心。较稳妥的历史判断是：君主保护一旦中断，已被触动的利益集团能迅速把改革者孤立；肃王此后对射中王尸者的处置，则显示新君同样不能放弃王室仪式的权威。
\end{debate}

吴起之死并没有使楚立刻崩解。疆域纵深使它还能吸收一次宫廷反转，北方防线也仍须维持。可是，这一结局留下了一条长期的反馈：每逢楚廷试图把封君资源、官吏责任和军粮更紧地拉回中心，就会面对谁来执行、谁会失去既得位置的问题。前394年前后调整俸禄与任免层次的尝试，\eventanchor{WQ-0394-CHU-DAO-PERSONNEL}{战国·楚悼王调整人事}，以及吴起改革前后所显出的反弹，不能证明所有改革从此消失，却足以说明楚的难处不在“不知道如何富国强兵”，而在于它难以让一次触动精英的重组越过君主更替而持续。

楚并非只是在北方承受压力。约前334年，楚威王北伐越国，越王无疆战败，越地原有的王权网络随之分裂，\eventanchor{WQ-0334-CHU-YUE-CONQUEST}{战国·楚越战争}；《史记·越王句践世家》又把无疆之死以及越人散入楚境、形成闽越等地方势力的后果连在一起，\eventanchor{WQ-0334-YUE-WUJIANG-DEATH}{战国·越王无疆败亡}。这场扩张使楚在淮泗与东南航路的影响一度扩大，却不能写成一次边界清楚、地方立即服从的吞并。《楚世家》与后世地理材料对于灭越时间和控制范围并不完全相同；所谓“分散”也不是一份可以逐户追踪的迁徙簿。较可把握的是，越不再维持统一王权，楚须在新近连结的东部地区面对不同的地方政治共同体。

\subsection{北方防线、合纵与秦楚之间的损耗}

吴起之后的楚仍是列国必须计算的力量。楚威王死后，前329年楚怀王即位，\eventanchor{WQ-0329-CHU-HUAI}{战国·楚怀王即位}；他面对的是齐、韩、魏在北方的牵制，和秦从西面不断施加的压力。前313年秦楚战事在丹阳、蓝田留下败楚的编年节点，\eventref{WQ-0312-QIN-CHU}{战国·丹阳蓝田之战}。丹阳失利后，怀王仍试图在前312年再攻蓝田，秦军反击使楚军东撤，\eventanchor{WQ-0312-CHU-LANTIAN-RETREAT}{战国·楚军退蓝田}；同一轮战事中屈匄等楚将被俘，\eventanchor{WQ-0312-CHU-QU-GAI-CAPTURE}{战国·丹阳俘将}。《楚世家》《秦本纪》与《资治通鉴》可以相互校出连续败北的脊柱，战场地望、战斗轮次和俘将姓名却仍有异文。它不宜被压缩为张仪一句游说造成的后果：《战国策》保存的说辞和人物言辞有强烈的修辞、编纂层次，不能逐字还原外交会谈。较能确认的是，秦以军事压力配合外交选择，楚在齐秦之间反复调整时，西部与汉水方向的缓冲已受到损耗。

这也解释了齐、韩、魏为何能在前303年至前301年协调行动。三国共同制楚，\eventanchor{WQ-0303-QI-HAN-WEI}{战国·齐韩魏制楚}；前301年联军在垂沙败楚，\eventanchor{WQ-0301-CHUISHA}{战国·垂沙之战}。传世材料还把楚将唐昧的被杀或被俘系于这场败战，\eventanchor{WQ-0301-CHU-TANGMEI-DEATH}{战国·唐昧失于垂沙}；《楚世家》与《战国策·楚策》对其结局和地名对应并不一致，因而只能把它视作楚廷失去北部防线统帅的一种强烈记忆，不能补成一份确切战报。联盟并非天然长久：三国各有对楚的顾虑，战后也没有形成稳定的共同政府或永久前线。它的作用在于楚已在对秦竞争中付出代价时，使其北方资源还要分给另一条战线。楚的广阔田地和城邑能够支撑恢复，却也意味着朝廷须同时供给方城、淮北、汉水和各条转运路线；一次多国协同，便足以放大这种调度负担。

怀王后期的危机由此不能只归咎于一次个人受骗。前299年，芈戎攻楚取新市，齐、魏、韩又攻方城，\eventanchor{WQ-0299-QIN-MIRONG-XINSHI}{战国·秦取新市}\eventanchor{WQ-0299-FANGCHENG-COALITION}{战国·方城受攻}；在土地争端与议和未解之际，怀王入秦。秦昭王将他扣留，并以割地施压，使一次出行转为人质危机，\eventanchor{WQ-0299-CHU-HUAI-DETAIN}{战国·秦扣留楚王}；怀王终于客死秦地，\eventanchor{WQ-0299-HUAIWANG}{战国·楚怀王被留秦地}。随从与楚廷在边境的交涉仍在继续，\eventanchor{WQ-0299-CHU-HUAI-RETINUE}{战国·楚廷交涉怀王}，但王位、救援与守边已不得不同时争夺资源。

《楚世家》《屈原列传》与《战国策》留下了丰富的劝说、怨愤和人物评断；它们可以说明后世如何把这场危机理解为外交判断的失败，却不能使我们听见秦楚廷中逐句说过什么。结构性原因是楚被置于多面压力之下，触发因素是怀王亲入秦使国君人身也成为谈判筹码，直接后果则是王权的可见性和对外号召力一并受损。

\subsection{失郢而不亡：东移的韧性与代价}

太子横在前298年即位为顷襄王，\eventanchor{WQ-0298-CHU-QINGXIANG-ASCEND}{战国·顷襄王即位}；楚廷以陈地为政治重心、重组宫廷和北部守军，见于\eventanchor{WQ-0299-CHU-CHEN-CAPITAL}{战国·楚廷转向陈地}与\eventanchor{WQ-0298-CHU-XIANG-ACCESSION-COURT}{战国·顷襄王重组朝廷}所概括的脉络。陈地的重心形成和后来失郢后的迁移，在传世材料中并非一条毫无空隙的行政记录；因此不宜画出王廷某日整体搬离、次日即在新地照旧运行的图景。可以确认的是，怀王危机之后，楚需要把官属、守军、仓储和王室的行动重心从受压的西北方向重新安排。

秦的推进使这项安排愈发急迫。前292年白起取宛，\eventanchor{WQ-0292-BAIQI-WAN}{战国·白起取宛}；《白起王翦列传》与南阳城址研究所能支持的，是秦把宛作为继续南进的据点、在当地留下守军，\eventanchor{WQ-0292-QIN-WAN-GARRISON}{战国·秦军驻宛}，而非留守规模与城内处置的完整清单。约前281年，白起又沿汉水方向推进，试探郢都以北的城链和水路补给，\eventanchor{WQ-0281-BAIQI-HANSHUI-ADVANCE}{战国·白起汉水推进}；前280年秦取黔中，\eventanchor{WQ-0280-QIANGZHONG}{战国·秦取黔中}，楚军从山地河谷退却，原先通向巴蜀、湘西的边防网络受阻，\eventanchor{WQ-0280-CHU-QIANZHONG-RETREAT}{战国·楚军退黔中}。撤军路线与地方反应没有连续记录，不能把它画成整齐的一条退线。楚廷则在陈地补置官属、仓储，\eventanchor{WQ-0280-CHU-CHEN-RELOCATION}{战国·楚廷补置陈地}；顷襄王朝继续向东南调动宗室、官属与物资，\eventanchor{WQ-0281-CHU-ROYAL-EXILE}{战国·楚廷东南调动}。这些节点不表示楚还有一张可以任意移动的资源清单；相反，旧郢及其西部资源越难直接动员，临时重建的宫廷、仓储和守备就越消耗人力、运输与地方合作。

\begin{event}{公元前278年以后}{郢都陷落与楚廷的东向重组}{可靠纪年，迁移过程有细节争议}{郢、陈地及江淮之间}
白起攻破郢都，\eventanchor{WQ-0278-YING}{战国·郢都失陷}。《史记·白起王翦列传》与纪南城考古研究共同提示，这一进攻冲击了楚王室的政治与礼仪中心，\eventanchor{WQ-0278-BAIQI-YING-TEMPLE}{战国·郢都宗庙受创}；焚毁范围和宗庙如何处置却没有同时代证据，不能据此罗列失去的祭器或府库。对任何以都城为行政和礼制中心的国家而言，这不只是失去一座城：宗庙、官属、王室成员、文书与可供调度的仓储若不能保全，政治共同体便可能随守城失败一同裂散。传世叙事所见楚廷转移这些核心，成为\eventanchor{WQ-0278-CHU-POST-YING-COURT}{战国·失郢后的楚廷}所要把握的事实；迁往陈地后安置王室祭祀、官属与仓储的过程，则由\eventanchor{WQ-0278-CHU-CHEN-RITUAL}{战国·楚廷迁陈}概括。陈城考古和《楚世家》无法提供迁移批次或礼器清单，却足以让人看见楚失去传统中心后仍要在淮北交通线上重建政令与祭祀的连续性；楚没有立即亡国。

秦随后取巫、黔中一带，\eventanchor{WQ-0277-QIN-CHU-WU}{战国·秦取巫黔中}。楚不得不更多使用江汉、淮水之间不经秦控制区的路线联结残余军政区域，\eventanchor{WQ-0277-CHU-SOUTHERN-ROUTES}{战国·楚改用南方路线}。南方的纵深因而确实延长了楚的生存，却把同样的难题带到更远处：一条能够避开敌军的水陆路线，也意味着王廷要在更多节点筹粮、验信、换船、联络封君和守将。东移是韧性，不是恢复旧日资源密度的同义词。
\end{event}

\begin{sourcenote}[郢都、陈地与“东迁”的材料边界]
\sourcebook{史记·楚世家}与《白起王翦列传》提供郢陷、楚廷延续及秦继续南进的编年骨架；《水经注》可帮助辨认后来叙述中的水陆空间，但成书更晚，不能替代战国行政档案。陈地何时成为稳定朝廷中心、官属仓储如何迁置、哪条路线由谁控制，均难逐项复原。本文只据材料能够支持的连续性，说明楚在失郢前后重新组织政治和运输中心，不把“东迁”写成一次没有摩擦的整体搬运。
\end{sourcenote}

\subsection{寿春的王廷与末年的选择}

失郢后的楚并非只会后退。前263年顷襄王卒，考烈王继位，\eventanchor{WQ-0263-CHUKAOLIE}{战国·楚考烈王即位}；《春申君列传》与《楚世家》将黄歇受任令尹、封为春申君，连到新王处理军政和对秦关系的需要，\eventanchor{WQ-0263-CHUNSHEN-PRIME-MINISTER}{战国·黄歇居令尹}。令尹任命、封地和早期政策的细节不尽相同，但春申君等相国网络确在新都、军队和使聘之间承担更重的协调工作，\eventanchor{WQ-0263-CHU-KAOLIE-COURT}{战国·春申君整合楚廷}。前257年，春申君遣楚军参与救赵，\eventanchor{WQ-0257-CHU-RESCUE-ZHAO}{战国·楚军救赵}。这不是楚已重获北方主导权的证明，而是秦若攻破赵国，压力会进一步东移、南压时，楚仍能把一部分资源送入合纵战争。

相国网络能把王室、军队和外交暂时接合，也会把协调权集中到宫廷少数人手中。前241年前后，考烈王由陈迁都寿春，\eventanchor{WQ-0241-CHU-SHOUCHUN}{战国·楚迁寿春}；寿春楚王城遗址报告与《楚世家》表明，新王畿须营建宫室、仓储和淮河舟运节点，\eventanchor{WQ-0241-CHU-SHOUCHUN-PALACE}{战国·寿春新王畿}。城址工程分期复杂，不能把所见遗存全定在迁都当年；围绕淮河的仓储、舟运与新王畿防务，仍成为\eventanchor{WQ-0240-CHU-SHOUCHUN-GRAIN}{战国·寿春的淮河防务}的条件。寿春缩短了王畿面临的防线，水网也有利于供给和调兵；代价是它离郢及楚的旧有政治资源更远。迁都后，春申君又联络赵、魏、韩、燕，楚军北上加入攻秦，\eventanchor{WQ-0241-CHUNSHEN-FIVE-STATES}{战国·春申君合纵}。五国并未形成由他统一统率的常备联军，实际承诺和统军程度尚待辨析；攻秦受挫更显示楚的远距离投送与联盟协调仍有限。迁都也没有消除封君、地方官和运输节点之间的协商成本，只是把成本重新分布到淮河沿线和末年的新中心。

前238年春申君被杀，\eventanchor{WQ-0238-CHUNSHEN-DEATH}{战国·春申君之死}；《春申君列传》与《楚世家》又将李园杀黄歇、扶持其所生子继为楚幽王串成一次相权与王位继承的重组，\eventanchor{WQ-0238-CHUNSHEN-CHULIE}{战国·李园杀春申君}。前228年楚幽王死后，负刍夺位，\eventanchor{WQ-0228-CHU-FUCHU-USURPATION}{战国·负刍夺位}，寿春宫廷的继承斗争随之削弱了应对秦军南进的准备，\eventanchor{WQ-0228-CHU-FUCHU-PALACE}{战国·负刍夺位后的寿春宫廷}。两事都不必被写成一位近臣或一场阴谋单独毁掉楚国的故事。《春申君列传》《楚世家》对于继承、李园和宫廷关系的铺陈，形成于事件之后，尤其人物动机、夺位过程和宫门中的言辞不可当作现场记录。更稳妥的是看它们所保存的政治后果：失郢后本已依靠相国网络维系的协调，在继承争夺中再次断裂；当秦军迫近时，王廷更难将淮南的资源、边将的判断和中央的战略接成连续行动。

\subsection{楚亡与“楚”没有消失}

前224年，李信伐楚受挫，\eventanchor{QIN-0224-CHU}{战国·李信伐楚受挫}。楚的纵深、动员能力以及项燕等抵抗力量，使秦不能把南征当作一次沿都城直线推进的收尾；《秦始皇本纪》与《白起王翦列传》的后出叙事又把昌平君在楚地反秦同这次反击联结，\eventanchor{WQ-0224-LIXIN-CHANGPINGJUN}{战国·李信失利与昌平君传说}。有关昌平君被拥为楚王的说法，见于\eventanchor{WQ-0224-CHANGPINGJUN-CHU}{战国·昌平君反秦传说}，其身份、称王与战役的关系仍有争议，不能确认为单一因果。无论如何，秦随即改以王翦、蒙武重整兵力与粮运，次年俘楚王负刍，\eventref{QIN-0223-CHU}{秦·灭楚}。《项羽本纪》与《白起王翦列传》还流传项燕在楚亡之际战死或自杀的说法，\eventanchor{WQ-0223-XIANGYAN-DEATH}{战国·项燕战死传说}；其死因、地点及与昌平君的关系均有重大争议，但楚军最后的统帅网络至此已难维系。

\sourcebook{史记·白起王翦列传}中王翦请“六十万”的数字，和列传安排的君臣对答，适合提示秦廷认识到南征需要比原先更多的兵力、后勤和耐心；它不是可以直接当作军籍总数的统计。楚的终结也不是疆域一大便必然失败的反证。正是这片纵深使它在郢陷后继续存在，使李信受挫，并迫使秦付出更大动员；但长期丧失西部与中原支点、宫廷继承反复、交通与资源协调成本上升，使楚难把广阔的人地稳定地转化为可以持续对抗秦的能力。

秦的占领没有使“楚”这个政治记忆立即消散。前208年项梁拥立熊心为楚怀王，\eventanchor{QIN-0208-HUAIWANG}{秦末·再立楚怀王}；项梁败死定陶后，怀王重配宋义、项羽、刘邦等人的任务，\eventanchor{QIN-0208-DINGTAO-CHU-REORGANIZATION}{秦末·楚军重组}，并以北救赵与西入关的双线命令维系联盟，\eventanchor{QIN-0208-HUAIWANG-DUAL-ORDERS}{秦末·楚怀王分兵}。这不是战国楚国原封不动地复活，而是旧国名号、旧地人群和新的军事集团在反秦危机中相互借用。

安阳驻军、宋义与项羽之间围绕时机和军粮的争议，\eventanchor{QIN-0208-CHU-SONGYI-CAUTION}{秦末·宋义观望}，以及项羽渡漳后攻击秦军甬道与粮运的行动，\eventanchor{QIN-0207-XIANGYU-JULU-SUPPLY-ATTACK}{秦末·楚军断秦粮道}，都出自《项羽本纪》《高祖本纪》的强叙事段落。它们可以显示补给与指挥权如何成为联盟的核心，不能用后出的激昂言辞填成未经证实的军前对话。楚亡后的这些事件反而提示：秦所接收的并非一片没有记忆的南方土地；曾经难以被单一中心完全整合的广阔社会与交通网络，在帝国危机时也能为新的政治名义提供人、路与动员的可能。

楚威王北伐越、越王无疆败死，\eventanchor{WQ-0334-CHU-YUE-CONQUEST}{战国·楚伐越}\eventanchor{WQ-0334-YUE-WUJIANG-DEATH}{战国·越王无疆败亡}，是楚向江淮扩张的一次节点，却不能据此把越地写成已被楚廷整齐接管的省份。部分越人进入楚境，部分地方网络后来与闽越等名号相连，皆是更长的政治重组；《越王句践世家》与后出的地理材料并不足以给出一张即时的人口迁徙图。前312年丹阳、蓝田的败退中，屈匄被俘或失陷、楚军东撤，\eventanchor{WQ-0312-CHU-QU-GAI-CAPTURE}{战国·丹阳失将}\eventanchor{WQ-0312-CHU-LANTIAN-RETREAT}{战国·蓝田东撤}，才让楚的西线压力从一次交锋变成对将领、城链与道路的持续损耗。

垂沙之后唐昧的死或被俘，\eventanchor{WQ-0301-CHU-TANGMEI-DEATH}{战国·垂沙失将}，和怀王入秦被扣，\eventanchor{WQ-0299-CHU-HUAI-DETAIN}{战国·怀王被扣}，都由传世叙事赋予很强的人物色彩；可见的政治后果更实际：北方守备须重新安排，王廷还要在国君人身成为外交筹码时维持粮道与号令。白起取宛、沿汉水推进，\eventanchor{WQ-0292-QIN-WAN-GARRISON}{战国·宛城驻军}\eventanchor{WQ-0281-BAIQI-HANSHUI-ADVANCE}{战国·秦军沿汉水推进}，又使失郢前的楚不能只以一座都城的安危来计算战争。

郢陷后迁陈，王室礼器、官属和仓储如何安置，\eventanchor{WQ-0278-BAIQI-YING-TEMPLE}{战国·郢都宗庙受冲击}\eventanchor{WQ-0278-CHU-CHEN-RITUAL}{战国·陈地安置王廷}，是材料只能给出方向、不能列出清册的工作。前263年黄歇为令尹，前241年参与五国攻秦并面对寿春新王畿的舟运与仓储，\eventanchor{WQ-0263-CHUNSHEN-PRIME-MINISTER}{战国·春申君为令尹}\eventanchor{WQ-0241-CHUNSHEN-FIVE-STATES}{战国·春申君参与合纵}\eventanchor{WQ-0241-CHU-SHOUCHUN-PALACE}{战国·寿春营建}，说明相国网络在迁都后能暂时把外交、军队和地方资源接在一起。前238年黄歇被杀、前228年负刍夺位，\eventanchor{WQ-0238-CHUNSHEN-CHULIE}{战国·春申君遇害}\eventanchor{WQ-0228-CHU-FUCHU-PALACE}{战国·寿春继承危机}，则使这一脆弱的接合再度裂开。李信受挫后昌平君与楚军的关联、项燕的最后结局，\eventanchor{WQ-0224-LIXIN-CHANGPINGJUN}{战国·昌平君反秦传说}\eventanchor{WQ-0223-XIANGYAN-DEATH}{战国·项燕终局}，尤其不能越过传世材料的限度；它们最多说明秦灭楚不是一条无阻的南下道路。

\begin{turningpoint}[制度得失：大国为何仍会失去主动]
\term{所得}：楚的江汉、江淮与南方纵深使它能够吸收宫廷政变、失郢和多次战败；玺印、简牍、仓储、渡口与地方官又表明，楚并非没有把资源转成命令的工具。吴起改革试图把封君、官吏与军费重新连到王权，失郢后的陈地、寿春与水路重组则一次次延长了国家的生命。

\term{代价与边界}：资源越分散，中心越需要在封君、官吏、编户与运输节点之间不断谈定执行；每次清核、征粮、迁置和守备，成本都落在具体地方关系上。吴起的结局显示改革若紧系一位君主便难越过继承，春申君之死与负刍夺位又显示末年协调无法自动延续。楚的疆域没有使它弱小；它给了楚抵抗的时间，却不能保证这些时间会凝结为更稳定的财政、军令和联盟。秦最终面对的正是一个仍有纵深的对手，楚最终失去的则是将纵深持续组织为国家能力的条件。
\end{turningpoint}
